% BHU PYQ -- Auto-generated & error-corrected
\documentclass[11pt,a4paper]{article}

\usepackage[a4paper,top=2.5cm,bottom=2.5cm,left=2.8cm,right=2.8cm,headheight=14pt]{geometry}
\usepackage{amsmath,amssymb}
\usepackage[T1]{fontenc}
\usepackage[utf8]{inputenc}
\usepackage{lmodern}
\usepackage{microtype}
\usepackage{fancyhdr}
\usepackage{enumitem}
\usepackage{hyperref}
\usepackage{lastpage}
\usepackage{parskip}

\hypersetup{colorlinks=true,urlcolor=black,linkcolor=black,
  pdftitle={Paper-VIII: Numerical Analysis and Computer Programming (Old Course) -- BHU PYQ},pdfauthor={Banaras Hindu University}}

\pagestyle{fancy}\fancyhf{}
\renewcommand{\headrulewidth}{0.4pt}\renewcommand{\footrulewidth}{0.4pt}
\fancyhead[L]{\small\textit{Banaras Hindu University}}
\fancyhead[R]{\small\textit{Paper-VIII: Numerical Analysis and Computer Pro}}
\fancyfoot[C]{\small Page \thepage\ of \pageref{LastPage}}

\newlist{parts}{enumerate}{2}
\setlist[parts,1]{label=(\alph*),leftmargin=*,itemsep=5pt,topsep=3pt}
\setlist[parts,2]{label=(\roman*),leftmargin=*,itemsep=3pt,topsep=2pt}
\newcommand{\pts}[1]{\hfill\textit{[\#1]}}

\begin{document}

\begin{center}
  {\large\bfseries BANARAS HINDU UNIVERSITY}\\[2pt]
  {\normalsize B.A./B.Sc. (Hons.) Part-III Examination, 2014-15}\\[6pt]
  \rule{0.8\linewidth}{0.5pt}\\[6pt]
  {\large\bfseries Paper No. Paper-VIII: Numerical Analysis and Computer Programming (Old Course)}\\[4pt]
  {\normalsize Time: 3 Hours\hfill Maximum Marks: 70}
\end{center}
\rule{\linewidth}{0.4pt}
\smallskip
\noindent\textit{Instructions: Answer \textbf{five} questions including Question~1 which is compulsory. Symbols have their usual meanings.}
\bigskip

BT/495
\begin{parts}
  \item Derive the Newton's forward difference interval formula for equal intervals. TAM see \& fd Sea S as saad wa at efter asia Examination, 2014-15 . Statistics Unit - II Sang -T Old Course Paper : VIII 4. Derive the general quadrature formula and
hence obtain Simpson's one third rule from it.
frat  vep-foes faq ast ort aise
Numerical Methods 5. Establish Newton-Cotes formula for numerical receipt of this question paper)
integration and herce obtain Trapezoidal rule A Note: Answer five questions in all. Question.No.
from it. 1 is compulsory. Answer ONE question from
aif Carat \& fra Best Hea Gs Bl waa each unit. All questions carry equal marks.
asa daw saa sitsigssra Faq at ora apis
Beat wa sett \& sere dis we Ge \pts{1}
ama 2\&1 7a sae A UH WS BH GAIT
Unit - ITI . i :
Be aA ust \& sip sar \pts{1}
gpg - Ill 6. Sum the following series to n terms :
  \item ve oa and extrapolation. aan AP Bon va ar ast sa wT: SF at a afedsr at orearer difere |
  \item Find the relationship between a and D
\end{parts}
\medskip\rule{\linewidth}{0.2pt}

\begin{parts}
  \item 5, 10, 17, 28, 47, 82, ......., operator. aD ort \& de aera ore aha BT/495
  \item Define factorial polynomial. (d (e (g (h
  \item  ) my Ua gese at oars diss | What is inverse interpolation ?Qt aa wT Ss ?Write the expression for the difference of product of two functions. a ort S WM H seat BH fra Shorey ferry | Evaluate 393 BT AT ATT HUT | Evaluate 2ab When interval of differencing is 1. SIT Dl SHS AMAHLC 2gh Hl AI aT ase | Find the general term of the series. Pet seh pr age ve att Hise
\end{parts}
\medskip\rule{\linewidth}{0.2pt}

2+ 6.4 +12.8+ 20.16 + 30.32+.....

Solve the following difference equation.
Fa oetetinct a ea aia
Hoa - Lux =0; x>0
x

BT/495
\begin{parts}
  \item Distinguish between the differential equations and difference equations. apa wich gq acdtipeit \& dra Ae aaa | Unit - I go - I
\end{parts}
\medskip\rule{\linewidth}{0.2pt}

\begin{parts}
  \item Ifinterval of differencing is one, prove that Shs Cee \& faa fae aia fe 'Of Bah) LIBS nan Sigg 'Most aM sd eth tency : Fafeted ada at fax aps :
\end{parts}
\medskip\rule{\linewidth}{0.2pt}

\begin{parts}
  \item Show that divided differences are symmetric functions of their arguments. aaa fe fasted sae oto Prafarstt \& aah Get ad S1 9.
  \item Describe Picard's method for solving ordinary differential equations. AIT Aaa Gti pw sq SA H fea fiers fafa ar arm aia Obtain the expression for remainder term in lagrange's interpolation formula. wast oddest dat A eT ve H fel ap OTT ast |
\end{parts}
\medskip\rule{\linewidth}{0.2pt}

200

\begin{parts}
  \item 
  \item  BT/495
\end{parts}
\medskip\rule{\linewidth}{0.2pt}

5 \pts{7}

+ + +...
25.8 5.8.11 8.11.14
Solve the following difference equations.
Pafeted aardhect sr ea apis
(
Meo tHoi tH, = x 2"

(ii)
(iii)
2 \pts{2}

Unit - IV
ge - IV

\begin{parts}
  \item Obtain a cube root of 10 using Newton Raphson method. ea UH false art 10 a Fea OTT aitsre | Discuss Gauss-Jordan method for solving simultaneous linear equations. ard Vas cect a et HA H fe Tea-stsa fate ar aia aise |
\end{parts}
\medskip\rule{\linewidth}{0.2pt}

\vfill
\rule{\linewidth}{0.4pt}
\begin{center}\small
  \href{https://bhu-exam.vercel.app/}{Click here to visit our website}\\[4pt]\href{https://me-aryan.vercel.app/}{Click here to visit our contributor}\\[4pt]\href{https://quashan-me.vercel.app/}{Click here to visit our contributor}
\end{center}

\end{document}